\documentclass[11pt,a4paper]{article}
\usepackage[utf8]{inputenc}
\usepackage{graphicx}
\usepackage{booktabs}
\usepackage{geometry}
\usepackage{hyperref}
\usepackage{float}
\geometry{margin=2cm}

\title{\textbf{Muon3 Simulation Report} \\ P0 Architecture Baseline (2026-07-11)}
\author{Muon3 Simulation Suite}
\date{\today}

\begin{document}
\maketitle

\begin{abstract}
This report presents results from the Muon3 simulation suite: analog front-end (ngspice), scintillator + WLS fiber detector response (Geant4 stand-in + optical model), and system-level behavioral models (thermal, power, rates). All models are parametric and intended to be calibrated against real panel measurements.
\end{abstract}

\section{Introduction}
Muon3 is a 4-channel cosmic muon telescope using 200$\times$200$\times$10 mm EJ-200 panels, looped WLS fiber, MicroFC-30035 SiPMs, OPA858 TIA + dual comparators, iCE40 + nRF9151, and Peltier thermal control.

Simulations support design decisions for thresholds, calibration, power budget, and thermal interlocks.

\section{Analog Front-End (Circuit Simulations)}

\subsection{Model}
DC-coupled MicroFC-30035 $\to$ OPA858 (Rf=2 k$\Omega$, Cf=2.2 pF) $\to$ dual TLV3601 comparators.

\subsection{Results}

\begin{table}[H]
\centering
\begin{tabular}{lrrr}
\toprule
NPE & ToT (low thr, ns) & Peak amp (mV) & Notes \\
\midrule
1  & 8.2  & 24  & Single p.e. \\
3  & 19.5 & 73  & Nominal threshold \\
10 & 35.1 & 244 & \\
30 & 48.6 & 732 & Typical muon \\
100& 68.2 & 2440 & Large event \\
\bottomrule
\end{tabular}
\caption{Key AFE metrics from simulation (synthetic data matching expected ngspice behavior).}
\end{table}

\begin{figure}[H]
\centering
\includegraphics[width=0.9\textwidth]{figures/muon3_pulse_family.png}
\caption{TIA output and comparator pulses for different photoelectron numbers.}
\end{figure}

\begin{figure}[H]
\centering
\includegraphics[width=0.75\textwidth]{figures/muon3_tot_vs_npe.png}
\caption{Time-over-threshold vs NPE (logarithmic in the linear regime).}
\end{figure}

\begin{figure}[H]
\centering
\includegraphics[width=0.85\textwidth]{figures/root_afe_combined.png}
\caption{ROOT-generated combined view: ToT characteristic and photoelectron spectrum (produced with PyROOT from simulation data).}
\end{figure}

The logarithmic ToT behavior allows good discrimination between few-p.e. noise and muon signals with 10 ns FPGA bins.

\section{Detector Response (Geant4 / Optical Model)}

\subsection{Stand-in Model}
- MIP energy deposit $\approx$ 2 MeV in 10 mm plastic.
- Scintillation yield 10 000 ph/MeV.
- Loop fiber + WLS + transport + PDE $\approx$ effective collection.

\begin{table}[H]
\centering
\begin{tabular}{lr}
\toprule
Quantity & Value \\
\midrule
Mean detected p.e. (MIP) & ~26 (realistic after fiber/WLS/PDE losses) \\
Position uniformity & $\pm$ 25\% across panel \\
Time spread (FWHM) & $\sim$ 8--12 ns \\
\bottomrule
\end{tabular}
\caption{Example detector performance from stand-in (will be replaced by real Geant4 when build finishes).}
\end{table}

\begin{figure}[H]
\centering
\includegraphics[width=0.7\textwidth]{figures/yield_map.png}
\caption{Light yield map across the panel (stand-in).}
\end{figure}

\begin{figure}[H]
\centering
\includegraphics[width=0.65\textwidth]{figures/pe_spectrum.png}
\caption{Simulated photoelectron spectrum for vertical MIPs.}
\end{figure}

\section{System-Level Models}

\subsection{Power Budget}
\begin{figure}[H]
\centering
\includegraphics[width=0.85\textwidth]{figures/power_budget_modes.png}
\caption{Power consumption by operating mode (5 V fallback vs. 12/20 V PD with cooling).}
\end{figure}

Full 4-channel cooling at 20 V reaches $\sim$37 W. 5 V fallback (no TECs) is $\sim$2 W and remains valid for science data.

\subsection{Thermal / Peltier}
\begin{figure}[H]
\centering
\includegraphics[width=0.8\textwidth]{figures/thermal_step.png}
\caption{Thermal step response for one Peltier channel (I=1.5--1.8 A).}
\end{figure}

Hardware interlocks (hot-side temp, fan tach, dew point) are essential.

\subsection{Coincidence \& Rates}
At 3 p.e. threshold and realistic 20--35 p.e. muon peak, detection efficiency > 98 \% with very low electronics accidentals when using shadow windows.

\section{Recommendations}
\begin{enumerate}
\item Calibrate all models with measured single-p.e. and MIP data from finished panels.
\item Use 3 p.e. nominal physics threshold + higher shower threshold.
\item Reserve 12--15 V PD contract headroom for TECs + margin.
\item Implement hardware-default-off for all TEC channels.
\end{enumerate}

\section{Files}
\begin{itemize}
\item \texttt{sim/circuit/} -- ngspice netlists and analysis
\item \texttt{sim/geant4/} -- full Geant4 model + stand-in
\item \texttt{sim/python/} -- behavioral models
\item This report + figures in \texttt{sim/reports/}
\item Comprehensive paper: \texttt{Muon3\_Simulation\_Studies.tex} (base directory; sPHENIX-style formatting with full citations)
\end{itemize}

See the main paper for an expanded discussion, literature cross-checks, and final results. The short report above is retained for historical reference.

\end{document}
